\par\clearpage
\noindent\makebox[0pt][l]{\fontsize{13}{16}\selectfont 416}\hfill{\small\textbf{POSTHUMOUS PAPERS.}}\hfill\null\par\medskip
\noindent tainly be important to subdivide the large southern triangles into
several smaller ones and to obtain a sound and sufficient basis for those
surveys by making measurements at new intermediate stations. I could make no
allowance for the costs of \textls[120]{these} operations in the above estimate. Estimating
them would be, if anything, more difficult still, but in any case they can only be
small compared with the costs of new large triangulations in the parts of the
country still untouched.

I need hardly remark that in future it will be desirable to make the results
for the positions of all the points precisely determined \textls[120]{public},
once they form a coherent whole. This goes without saying in any case for the
strictly scientific part of the measurements made so far.\quad\mbox{.\ .\ .\ .\ .\ .}

\bigskip
\begin{center}
[From the report of 26~June 1828, »die Fortsetzung der Gradmessungs-\par
arbeiten betreffend.«]
\end{center}

\par\noindent\hspace*{\parindent}\mbox{.\ .\ .\ .\ .\ .}\quad Since the detailed survey of the parts of the country still
unsurveyed is to be based on the trigonometrical operations, and the two
undertakings will depend on one another with regard to the funds to be used,
a rough estimate of the total costs was needed first. According to a note
communicated to me by Privy Cabinet Councillor \textsc{Hoppenstedt}, the area
of the parts of the country to be surveyed in detail would amount to about
144 square miles. The costs of the detailed survey by General Staff officers
are estimated at 200---250 thalers for each square mile, to which about 21
thalers would have to be added for the cost of making 4 copies of the map.
Thus, if the total were reckoned at 250 thalers per square mile, these costs
would amount to about 36000 thalers and consequently to nearly 50000 thalers
including the costs of triangulation in the parts of the country still untouched
by the arc measurement. One or a couple of thousand thalers might have to be
added for the operations needed to establish, within the large southern
triangles of the arc measurement, a sufficiently large number
\par
\clearpage
\noindent\null\hfill{\small\textbf{ON THE HANOVER TRIANGULATION.}}\hfill\makebox[0pt][r]{\fontsize{13}{16}\selectfont 417}\par\medskip
\noindent of fixed points for the detailed survey, a matter on which I already
gave a fuller account in the submission lodged in November of last year. Thus,
if 5000 thalers can actually be spent on these undertakings annually (which will
depend on the continuous availability of the personnel as well as on the funds),
a period of about 10 years would be needed to complete the work in full.

According to an earlier estimate of my own, though perhaps a less reliable one,
the area of the parts of the country to be surveyed in detail would be 173
square miles. On that basis the estimated cost would be about 7000 thalers
higher and the work would take 1 to 2 years longer. Likewise, if the allotted
5000 thalers are to be taken solely as the maximum for each year, with no
assumption that they are actually to be spent every year or on average, the
duration of the whole undertaking would be proportionately longer.

As regards the distribution of the trigonometrical work over the individual
years, to complete it more quickly as the basis for the detailed survey it
might be advisable at first to spend the larger part of the funds on it and the
smaller on the detailed survey, perhaps in the ratio of \GaussIXMathCopy{2/3}{\(\frac{2}{3}\)} to
\GaussIXMathCopy{1/3}{\(\frac{1}{3}\)}. The latter would then be assured of all the more rapid progress
in later years, when a sound basis is available everywhere, when the workers
have gradually gained more and more practice, and when the funds can be spent
on it alone. As regards the order in which the trigonometrical work is to be
carried out, as things stand there is no further reason to prefer one operation
to another beyond the need to ensure that the concurrent detailed survey is
always supplied with a sufficient stock of reliably determined points so that
it is never in danger of coming to a standstill for want of them. A more
specific determination might well be, \textls[120]{for the present}, in some
respects impracticable, in others unnecessary, and in still others not even
advisable, since a certain degree of freedom to work on whatever is most
appropriate in the circumstances at any given time can only promote faster
and better progress.
\par\noindent\mbox{.\ .\ .\ .\ .\ .}\par
\medskip\noindent\hspace*{0.03\linewidth}{\footnotesize IX.}\hfill\makebox[0.18\linewidth][l]{53}\par
